\documentclass[11pt]{article}
\usepackage[margin=1in]{geometry}
\usepackage{amsmath,amssymb,amsthm,mathtools,mathrsfs,booktabs,array}
\usepackage{microtype,xurl}
\usepackage{xcolor}
\usepackage[colorlinks=true,linkcolor=blue!55!black,citecolor=blue!55!black,urlcolor=blue!55!black]{hyperref}

\newtheorem{theorem}{Theorem}[section]
\newtheorem{proposition}[theorem]{Proposition}
\newtheorem{corollary}[theorem]{Corollary}
\theoremstyle{remark}
\newtheorem{remark}[theorem]{Remark}
\newcommand{\R}{\mathbb R}
\newcommand{\C}{\mathbb C}
\newcommand{\cT}{\mathcal T}
\newcommand{\cQ}{\mathcal Q}
\newcommand{\cO}{\mathcal O}
\newcommand{\rank}{\operatorname{rank}}
\newcommand{\Sol}{\operatorname{Sol}}
\newcommand{\Diff}{\operatorname{Diff}}
\newcommand{\vol}{\operatorname{vol}}
\newcommand{\cert}[1]{\nolinkurl{#1}}

\title{Einstein--Maxwell Waves inside Weyl--Maxwell Gravity\\
on a Compact Product:\\
Exact Linear Phase-Space Inclusion and the Extra Axial Branch}
\author{GPT-5\thanks{OpenAI model. Contributed the mathematical synthesis,
claim audit, and manuscript. The project was commissioned and directed by
Asger Alstrup Palm (\texttt{asger@area9.dk}), who initiated the questions,
coordinated the programme, and serves as corresponding human contact.}}
\date{Working draft, July 2026}

\begin{document}
\maketitle

\begin{abstract}
We locate ordinary Einstein--Maxwell gravitational waves inside pure
Weyl--Maxwell gravity on the compact product
$\R_t\times S^1_L\times S^2$. The background is a common exact solution of
both theories: the sphere has unit curvature, the Maxwell field is a unit
parallel magnetic flux, and the action normalizations are
$\kappa=1$, $\Lambda=1/2$, and $\alpha_B=3$. On the fixed compact $U(1)$
bundle we classify the complete standard linearized Einstein--Maxwell
harmonic phase space, including all axial and polar radiative modes, the
physical exceptional $\ell=1$ quotient, the six-dimensional homogeneous
generalized block, and the three axial twist pairs.

The identity tangent map into Weyl--Maxwell theory is injective and its
Lee--Wald pullback is nondegenerate on every standard block. It is not a
symplectic inclusion. For $\ell\geq2$ the target form is the
Einstein--Maxwell form composed with the parity-independent spectral
polynomial
\[
 p_\lambda(M)=1+\frac32(M-\lambda),
\]
whose two normal-mode weights are $1\pm\frac32\sqrt{2\lambda}$. Thus each
real spatial harmonic has relative signature $(2,2)$. The physical $\ell=1$
block instead acquires the factor $4$, the homogeneous block a nontrivial
unipotent shear, and every axial twist pair the factor $-2$.

We then solve the complete generic axial Weyl--Maxwell target. Modulo the
Einstein image, it contains two additional cyclic summands on
$\omega^2-k^2=\lambda-2/3$. Their direct four-dimensional Lee--Wald Gram
matrix is nondegenerate and positive in the declared positive-frequency
current convention. The Einstein and extra modules are mutually orthogonal,
and the complete generic axial target has classical signature $(3,1)$.
Inverting the extra Gram matrix gives two conserved coefficient detectors
which annihilate the Einstein image.

These are classical \textsc{local-algebraic}/\textsc{reduced-mode}
statements before the final residual quotient. They prove neither a
positive-frequency Hilbert space nor a quantum ghost theorem. They also do
not construct asymptotically flat scattering. The ordinary waves are present
and nonnull before the global residual reduction; a vanishing one-particle
residual cohomology on a closed cylinder therefore answers a different
question from the existence of local gravitational radiation.
\end{abstract}

\section*{Authorship and computational inputs}
The named model author synthesized the exact results and wrote the
manuscript. The mathematical inputs are independently generated exact
certificates in the project repository. Section~\ref{sec:repro} gives their
content hashes and verification commands. No floating-point calculation is
used in a rank, signature, quotient, or on-shell identity reported here.

\tableofcontents

\section{The question and the answer}\label{sec:intro}

Every four-dimensional Einstein metric is Bach-flat, but the converse is
false: conformal gravity has genuine non-Einstein solution branches
\cite{Bach,Maldacena,DeserJoungWaldron}. In the
coupled theory the relation is subtler still, because the Maxwell stress
appears in both equations with different gravitational operators. The
relevant question is consequently not whether one can fix Weyl gauge to
obtain Einstein gravity. One cannot do so in general. It is instead:
\begin{quote}
On a common Einstein--Maxwell/Weyl--Maxwell background, which ordinary wave
solutions embed, what phase-space form do they inherit, and which additional
Weyl--Maxwell solutions remain?
\end{quote}

This paper answers that question completely for the standard harmonic
Einstein--Maxwell tangent and for the generic axial part of the additional
target on a fixed compact magnetic bundle. The answer has three logically
separate layers:
\begin{enumerate}
\item every standard Einstein--Maxwell tangent class injects as a
Weyl--Maxwell solution class;
\item the target Lee--Wald form restricts nondegenerately to that image, but
does not equal the source form under the identity map; and
\item the generic axial Weyl--Maxwell target contains two further nonradical
solution directions which are not Einstein representatives or local gauge.
\end{enumerate}

Let $\cT_{\rm EM}^{\rm std}$ denote the complete fixed-bundle standard
Einstein--Maxwell harmonic tangent after local $\Diff\times U(1)$ reduction,
and let $\cT_{\rm WM}^{\rm ax}$ denote the generic axial Weyl--Maxwell tangent
after local $\Diff\times\text{Weyl}\times U(1)$ reduction. The central result
can be stated schematically as
\begin{equation}\label{eq:central-summary}
 \boxed{\begin{gathered}
 \iota:\cT_{\rm EM}^{\rm std}\hookrightarrow\cT_{\rm WM},
 \qquad \ker(\iota^*\Omega_{\rm WM})=0,\\
 \iota^*\Omega_{\rm WM}\neq\Omega_{\rm EM},\\
 \cT_{\rm WM}^{\rm ax}/\iota(\cT_{\rm EM}^{\rm ax})
 \cong \bigl(F[\omega]/(p)\bigr)^2,
 \qquad p=\omega^2-k^2-\lambda+\frac23 .
 \end{gathered}}
\end{equation}
Here $F=\operatorname{Frac}\R(\lambda,k)$ on the generic stratum. The last
line is a canonical module quotient, not a choice of complement.

\subsection{What is and is not frozen}
The theorem level of this manuscript is deliberately asymmetric in parity.
The standard Einstein--Maxwell image is complete in both axial and polar
parities and includes all exceptional/global blocks. The additional
Weyl--Maxwell branch is complete only in the generic axial sector. The polar
extra branch and the final residual descent remain open. Accordingly the
paper makes no parity-complete claim about the full Weyl--Maxwell target.

The dependency labels are
\begin{center}
\textsc{local-algebraic}\quad and\quad \textsc{reduced-mode}.
\end{center}
No result below carries a \textsc{lorentzian-causal} dependency label. The
compact background is Lorentzian and the currents are conserved in real
time, but no full Green-hyperbolic BV complex, Hadamard state, null-infinity
phase space, or scattering construction is used.

\section{The two theories and their common compact background}
\label{sec:background}

We use signature $(-,+,+,+)$ and the actions
\begin{align}
 S_{\rm EM}[g,A]
 &=\int_M\!\sqrt{-g}\left[
 \frac{R-2\Lambda}{2\kappa}-\frac14F_{ab}F^{ab}\right]d^4x,
 \label{eq:EM-action}\\
 S_{\rm WM}[g,A]
 &=\int_M\!\sqrt{-g}\left[
 \frac{\alpha_B}{8}C_{abcd}C^{abcd}-\frac14F_{ab}F^{ab}\right]d^4x .
 \label{eq:WM-action}
\end{align}
Thus
\begin{align}
 G_{ab}+\Lambda g_{ab}&=\kappa T_{ab},
 &\nabla_aF^{ab}&=0,\label{eq:EM-eom}\\
 \alpha_B B_{ab}&=T_{ab},
 &\nabla_aF^{ab}&=0,\label{eq:WM-eom}
\end{align}
with $dF=0$ and
\begin{equation}
 T_{ab}=F_{ac}F_b{}^c-\frac14g_{ab}F_{cd}F^{cd}.
\end{equation}

Consider a product $M_2(k_1)\times\Sigma_2(k_2)$ of oriented constant
curvature surfaces and an aligned parallel field
\begin{equation}
 F=E\,\vol(M_2)+P\,\vol(\Sigma_2).
\end{equation}
The same metric and field solve both theories on the incidence branch
\begin{equation}\label{eq:incidence}
 \Lambda=\frac{k_1+k_2}{2},\qquad
 E^2+P^2=\frac{k_2-k_1}{\kappa},\qquad
 \alpha_B\kappa(k_1+k_2)=3.
\end{equation}
This is an intersection of solution loci, not a gauge equivalence.

We specialize to the positive flat branch
\begin{equation}\label{eq:fixture}
 M=\R_t\times S^1_L\times S^2,\qquad
 (k_1,k_2)=(0,1),\qquad
 (\kappa,\Lambda,\alpha_B,E,P)=(1,\tfrac12,3,0,1).
\end{equation}
The spatial Cauchy surface is the closed manifold $S^1\times S^2$. With
minimal charge $q_{\min}=1$, flux quantization gives Chern number $N=2$.
Throughout the paper the bundle $P_N$ is fixed. Uniform magnetic variation
would change the topological sector and is not an allowed tangent.

\subsection{Why the complete linear inclusion holds}
For an Einstein--Maxwell solution, the difference between the Bach equation
and the Einstein equation can be expressed through the Chevreton tensor of
the Maxwell field. On the product fixture the background flux is parallel.
The Chevreton defect \cite{Chevreton} is quadratic in $\nabla F$, so both its background value
and first variation vanish. With the coupling relation \eqref{eq:incidence},
every complete linearized Einstein--Maxwell solution therefore satisfies the
complete linearized Weyl--Maxwell equations with the same pair $(h,a)$.

This is an on-shell tangent statement. It is stronger than principal-symbol
inclusion and weaker than an off-shell BV chain map. It also predicts its own
limitation: the first possible Chevreton failure is quadratic, so no
nonlinear closure follows from the linear theorem.

Let
\begin{equation}
 \iota[(h,a)]_{\Diff\times U(1)}
 =[(h,a)]_{\Diff\times\mathrm{Weyl}\times U(1)}.
\end{equation}
The map is well defined on the declared linear solution quotients. It is also
injective. Indeed, if an Einstein class becomes target gauge, then after
subtracting its common diffeomorphism and $U(1)$ part it has the form
$(h_{ab},a_a)=(2\sigma\bar g_{ab},0)$. The linearized Einstein--Maxwell rows
imply
\begin{equation}\label{eq:weyl-injectivity}
 -3\Delta_{S^2}\sigma+2\sigma=0.
\end{equation}
On $Y_{\ell m}$ this is multiplication by $3\ell(\ell+1)+2$, hence smoothness
forces $\sigma=0$.

\begin{proposition}[Linear quotient inclusion]\label{prop:inclusion}
On the fixed background \eqref{eq:fixture}, the identity tangent map sends
every complete linearized Einstein--Maxwell solution to a linearized
Weyl--Maxwell solution and descends injectively through the respective local
gauge quotients.
\end{proposition}

\begin{remark}
Proposition~\ref{prop:inclusion} does not say that the two theories are gauge
equivalent. In particular, a generic Bach-flat solution need not be
conformally Einstein, and Section~\ref{sec:extra} constructs additional
target directions explicitly.
\end{remark}

\section{The complete standard Einstein--Maxwell harmonic phase space}
\label{sec:source}

Write
\begin{equation}
 \lambda=\ell(\ell+1),\qquad k=k_n=\frac{2\pi n}{L},
 \qquad -\Delta_{S^2}Y_{\ell m}=\lambda Y_{\ell m},
\end{equation}
and $N_{\ell m}=\int_{S^2}|Y_{\ell m}|^2d\Omega>0$. A complex coefficient at
$(n,\ell,m)$ is paired by reality with the coefficient at $(-n,\ell,-m)$.
Equivalently one may use real spherical harmonics and cosine/sine circle
modes. We use the latter convention for real mode counts.

\subsection{Regular radiative blocks}
For $\ell\geq2$, the axial and polar reductions each have two masters. In the
orders $(H,Q)$ and $(K,U)$ their normal-mode matrices are
\begin{equation}\label{eq:master-matrices}
 M_A=\begin{pmatrix}\lambda&2\\ \lambda&\lambda\end{pmatrix},
 \qquad
 M_P=\begin{pmatrix}\lambda&-2\lambda\\ -1&\lambda\end{pmatrix}.
\end{equation}
Both parities have the exact dispersions
\begin{equation}\label{eq:dispersion}
 \omega_\pm^2=k^2+\lambda\pm\sqrt{2\lambda}.
\end{equation}
The source action-normalized Wronskians are controlled, up to the common
positive factor $N_{\ell m}/2$, by
\begin{equation}\label{eq:source-forms}
 G_A=\begin{pmatrix}\lambda^2&2\lambda\\2\lambda&2\lambda\end{pmatrix},
 \qquad
 G_P=\begin{pmatrix}1&-2\\-2&2\lambda\end{pmatrix}.
\end{equation}
They obey $G_AM_A=M_A^TG_A$ and $G_PM_P=M_P^TG_P$. Their determinants are
$2\lambda^2(\lambda-2)$ and $2(\lambda-2)$, so both are positive for
$\lambda\geq6$. This is positivity of the classical source kinetic branch
weights, not yet a one-particle Hilbert norm.

The branch eigenvectors may be taken as
\begin{equation}
 v_{A,\pm}=(1,\ \pm\sqrt{\lambda/2}),\qquad
 v_{P,\pm}=(\mp\sqrt{2\lambda},\ 1).
\end{equation}
These compact parity masters are the exact representation of the ordinary
coupled photon/graviton-like waves. On the magnetically supported product
the metric and Maxwell coefficients mix; helicity $\pm2$ is therefore a
local/high-frequency interpretation, whereas $(n,\ell,m)$, parity, and
branch are the exact global labels.

\subsection{Exceptional and global blocks}
At $\ell=1$ one branch becomes gauge in each parity. The axial source matrix
has kernel $(H,Q)=(1,-1)$ and physical direction $(1,1)$. The polar source
matrix has kernel $(K,U)=(2,1)$ and quotient coordinate
$\Psi_P=U-K/2$. The surviving axial and polar modes obey
$\omega^2=k^2+4$. They are radiative triplets, not the zero-frequency twist.

The complete homogeneous generalized block is six-dimensional. A convenient
representative is
\begin{equation}\label{eq:hom-rep}
 K=a+bt,\qquad C=at^2+\frac b3t^3+c+dt,\qquad A_x=W_x+Q_et.
\end{equation}
Polynomial Jordan partners are retained; imposing boundedness in time would
define a different phase space. After the common factor $2\pi L$ is removed,
the source form in the order $(a,b,c,d,Q_e,W_x)$ is
\begin{equation}\label{eq:hom-source}
 \Omega_{\rm EM}^{(0)}=
 \begin{pmatrix}
 0&-1&0&-1&0&0\\ 1&0&1&0&0&0\\
 0&-1&0&0&0&0\\ 1&0&0&0&0&0\\
 0&0&0&0&0&-2\\ 0&0&0&0&2&0
 \end{pmatrix}.
\end{equation}
It has rank six. The flat holonomy $W_x$ is the canonical partner of the
electric charge $Q_e$ and must not be deleted merely because it is invisible
to the field strength.

For every one of the three real axial $\ell=1$ harmonics, the generalized
twist
\begin{equation}\label{eq:twist-rep}
 h_{xa}=(A_m+B_mt)X_a,\qquad a_x=-(A_m+B_mt)Y_{1m}
\end{equation}
forms a Darboux pair. The would-be gauge generator is proportional to
$x(A_m+B_mt)$ and is not periodic on $S^1$.

\begin{theorem}[Complete source phase space]\label{thm:source-complete}
Before the final residual quotient, the fixed-bundle standard harmonic
Einstein--Maxwell tangent is the symplectic direct sum
\begin{equation}\label{eq:source-decomp}
 \cT_{\rm EM}^{\rm std}
 =\cT_{\rm rad}^{\ell\geq2}\oplus\cT_{\rm phys}^{\ell=1}
 \oplus\cT_{\rm hom}^{\ell=0}
 \oplus\cT_{\rm twist}^{\ell=1,\omega=0}.
\end{equation}
The regular radiative forms, exceptional physical quotient, homogeneous
block, and three twist pairs are all nondegenerate.
\end{theorem}

\section{The Weyl--Maxwell pullback on the Einstein image}
\label{sec:restriction}

A solution inclusion need not preserve a symplectic structure. We use the
covariant phase-space current in the Lee--Wald convention
\cite{LeeWald,IyerWald} and therefore
keep three objects separate:
\begin{equation}
 \Omega_{\rm EM},\qquad \iota^*\Omega_{\rm WM},\qquad
 \Omega_{\rm WM}\ \hbox{on the full target}.
\end{equation}
The first two live on the same vector space but arise from different actions.

\subsection{Regular radiation: one spectral polynomial}
On both parities and for every $\ell\geq2$,
\begin{equation}\label{eq:spectral-identity}
 \iota^*\Omega_{\rm WM}(u,v)
 =\Omega_{\rm EM}\bigl(u,p_\lambda(M)v\bigr),
 \qquad p_\lambda(x)=1+\frac32(x-\lambda).
\end{equation}
In the axial and polar master coordinates,
\begin{equation}
 p_\lambda(M_A)=
 \begin{pmatrix}1&3\\3\lambda/2&1\end{pmatrix},\qquad
 p_\lambda(M_P)=
 \begin{pmatrix}1&-3\lambda\\-3/2&1\end{pmatrix}.
\end{equation}
On the normal branches its eigenvalues are
\begin{equation}\label{eq:relative-weights}
 r_+=1+\frac32\sqrt{2\lambda}>0,\qquad
 r_-=1-\frac32\sqrt{2\lambda}<0.
\end{equation}
Neither vanishes for $\lambda\geq6$. The axial and polar branch weights are
identical even though their off-shell coefficient matrices differ.

For a real spatial harmonic there are four oscillator blocks: two parities
times two branches. Hence the relative coefficient signature is
\begin{equation}\label{eq:radiative-signature}
 \operatorname{sig}_{\rm rel}\bigl(\iota^*\Omega_{\rm WM}\bigr)=(2,2).
\end{equation}
This is shorthand for the signs of the symmetric branch coefficients that
multiply the standard real Darboux blocks. A symplectic form itself has no
ordinary Sylvester signature.

\subsection{Exceptional and global restrictions}
The exceptional $\ell=1$ calculation cannot be obtained by blindly setting
$\lambda=2$ in the generic polar matrix: that continuation fails target
gauge descent. Direct four-dimensional currents instead give, in
source-normalized axial and polar quotient coordinates,
\begin{equation}\label{eq:ell1-factor}
 \left.\iota^*\Omega_{\rm WM}\right|_{\ell=1,\rm phys}
 =4\left.\Omega_{\rm EM}\right|_{\ell=1,\rm phys}.
\end{equation}
The literal target current has zero gauge rows and columns, so this is a
quotient identity rather than only a norm comparison.

On the homogeneous block, after the same common factor as in
\eqref{eq:hom-source},
\begin{equation}\label{eq:hom-target}
 \Omega_{\rm WM}^{(0)}=
 \begin{pmatrix}
 0&2&0&-1&0&0\\ -2&0&1&0&0&0\\
 0&-1&0&0&0&0\\ 1&0&0&0&0&0\\
 0&0&0&0&0&-2\\ 0&0&0&0&2&0
 \end{pmatrix}.
\end{equation}
Writing $R=(\Omega_{\rm EM}^{(0)})^{-1}\Omega_{\rm WM}^{(0)}$, one obtains
\begin{equation}\label{eq:unipotent}
 R=I+N,\qquad N^2=0,\qquad\rank N=2.
\end{equation}
Thus the identity is not symplectic, although $S=I+N/2$ has determinant one
and satisfies $S^T\Omega_{\rm EM}^{(0)}S=\Omega_{\rm WM}^{(0)}$.

Finally, each generalized axial twist pair obeys the direct identity
\begin{equation}\label{eq:twist-factor}
 \left.\iota^*\Omega_{\rm WM}\right|_{\rm twist}
 =-2\left.\Omega_{\rm EM}\right|_{\rm twist}.
\end{equation}
It comes from the full $A+Bt$ current and is not a zero-frequency limit of an
oscillatory formula.

\subsection{Orthogonality and the assembled theorem}
Different circle momenta, spherical irreducibles, and parity sectors are
orthogonal by $S^1$ Fourier invariance, $SO(3)$ invariance, and spatial
parity. Distinct radiative branches are orthogonal because $M$ is
$\Omega_{\rm EM}$-self-adjoint and has distinct eigenvalues. The only shared
label not settled by these arguments is the physical axial
$\ell=1,n=0$ oscillator against the twist. Its literal mixed target current
is proportional to
\begin{equation}
 -2i\pi\omega p(\omega^2-4)
 [\omega(A+Bt)-iB]e^{-i\omega t},
\end{equation}
and therefore vanishes on the physical shell $\omega^2=4$ for both Jordan
coordinates.

\begin{theorem}[Complete standard harmonic inclusion]
\label{thm:standard-inclusion}
On the fixed compact background \eqref{eq:fixture}, and before the final
residual quotient, the identity map
\begin{equation}
 \iota:\cT_{\rm EM}^{\rm std}\hookrightarrow\cT_{\rm WM}
\end{equation}
has nondegenerate target pullback:
\begin{equation}
 \ker\left(\iota^*\Omega_{\rm WM}
       \big|_{\cT_{\rm EM}^{\rm std}}\right)=0.
\end{equation}
The pullback is the block-orthogonal direct sum of
\eqref{eq:spectral-identity}, \eqref{eq:ell1-factor},
\eqref{eq:hom-target}, and \eqref{eq:twist-factor}. The identity inclusion is
not symplectic.
\end{theorem}

\begin{corollary}[Restricted linear observables]
On the restricted Einstein image, every linear Einstein--Maxwell observable
has a unique target-restricted Hamiltonian vector after application of the
inverse blockwise relative operator.
\end{corollary}

\begin{remark}
The corollary is not an embedding into the full target observable algebra.
Extending an observable away from the Einstein image and descending it under
the final residual action are separate questions.
\end{remark}

\section{The complete generic axial extra branch}\label{sec:extra}

We now solve the fourth-order target rather than restricting it to the
Einstein image. For generic $\ell\geq2$, after local gauge contraction the
axial coefficient order is
\begin{equation}
 \Phi=(H_t,H_x,Q_t,Q_x)^T.
\end{equation}
The contraction inverts only the harmless constant $2$: neither $k$ nor
$\omega$ is inverted, so zero momentum is retained.

Put $s=\omega^2-k^2$. The Einstein master polynomial and the extra polynomial
are
\begin{equation}
 q(s)=(s-\lambda)^2-2\lambda,\qquad p(s)=s-\lambda+\frac23.
\end{equation}
Their resultant vanishes only at the algebraic collision $\lambda=2/9$,
which is disjoint from the physical range $\lambda\geq6$. The exact
Smith/module reduction gives
\begin{equation}\label{eq:extra-module}
 \cQ_{\rm ax}^{\rm extra}
 :=\cT_{\rm WM}^{\rm ax}/\iota(\cT_{\rm EM}^{\rm ax})
 \cong\bigl(F[\omega]/(p)\bigr)^2,\qquad
 F=\operatorname{Frac}\R(\lambda,k).
\end{equation}
Thus there are two independent additional cyclic summands on
\begin{equation}\label{eq:extra-shell}
 \omega_e^2=k^2+\lambda-\frac23.
\end{equation}
Convenient representatives are
\begin{equation}\label{eq:extra-reps}
 e_1=\begin{pmatrix}-k^2-\lambda\\k\omega_e\\\lambda\\0\end{pmatrix},
 \qquad
 e_2=\begin{pmatrix}-k\omega_e\\k^2-\frac23\\0\\\lambda\end{pmatrix}.
\end{equation}
They are not defined by an arbitrary orthogonal complement: they represent a
canonical quotient supported on a different invariant factor.

\subsection{Direct Lee--Wald pairing}
Normalize the direct current by $J^t/(-i\omega_eN_{\ell m})$ at positive
frequency, suppressing the common positive circle volume. In the basis
\eqref{eq:extra-reps} its Gram matrix $G_X$ has entries
\begin{align}
 (G_X)_{11}
 &=\frac{\lambda}{6}\bigl[
 9k^2\lambda^2-12k^2\lambda+4k^2
 +9\lambda^3-18\lambda^2\bigr],\label{eq:gx11}\\
 (G_X)_{12}
 &=\frac{k\lambda\omega_e}{6}(3\lambda-2)^2,\label{eq:gx12}\\
 (G_X)_{22}
 &=\frac{\lambda}{18}\bigl[
 27k^2\lambda^2-36k^2\lambda+12k^2
 +36\lambda-8\bigr].\label{eq:gx22}
\end{align}
On the extra shell,
\begin{equation}\label{eq:extra-det}
 \det G_X=\frac13\lambda^4(\lambda-2)(9\lambda-2)>0.
\end{equation}
The first principal minor is also positive for every real $k$ and
$\lambda\geq6$. Hence $\operatorname{sig}G_X=(2,0)$.

The Einstein and extra primary modules are mutually Lee--Wald orthogonal.
The two Einstein master branches contribute one positive and one negative
coefficient in the same direct-current convention. Therefore:

\begin{theorem}[Generic axial target classification]
\label{thm:axial-target}
For every $\ell\geq2$, every real compact momentum $k$, and before the final
residual quotient, the complete generic axial Weyl--Maxwell solution module
is the Lee--Wald orthogonal direct sum of the two Einstein master branches and
the rank-two extra module \eqref{eq:extra-module}. It is nondegenerate and has
classical coefficient signature
\begin{equation}
 \operatorname{sig}\cT_{\rm WM}^{\rm ax}=(3,1).
\end{equation}
The extra block has signature $(2,0)$; the negative target direction lies on
the minus Einstein-image branch.
\end{theorem}

\begin{remark}[What the sign does not prove]
Theorem~\ref{thm:axial-target} is a sign theorem for a classical conserved
current after a declared positive-frequency normalization. No compatible
complex structure, positive one-particle inner product, BRST physical Hilbert
space, or quantum time evolution has been constructed. Consequently the
signature $(3,1)$ is not, by itself, a quantum ghost or unitarity theorem.
The negative target coefficient belongs to $\iota^*\Omega_{\rm WM}$, not to
the independent Einstein--Maxwell source action. Familiar higher-derivative
ghost conclusions \cite{Stelle} require additional quantum structure not
constructed here.
\end{remark}

\subsection{A detector for the extra coefficients}
Nondegeneracy makes the extra branch measurable within the declared reduced
solution block. For an on-shell axial solution $\Phi$, define
\begin{equation}\label{eq:detector}
 \cO_X(\Phi)
 =G_X^{-1}\frac{\Omega_{\rm WM}(\overline e,\Phi)}
 {-i\omega_eN_{\ell m}}\in\C^2,
\end{equation}
where $e=(e_1,e_2)^T$. The local Green identity gives slice independence, and
direct substitution gives
\begin{equation}
 \cO_X(a_1e_1+a_2e_2)=(a_1,a_2)^T,\qquad \cO_X\circ\iota=0
\end{equation}
on the complete generic axial Einstein image. Thus the additional solutions
are neither a reduced-current radical nor a relabeling of Einstein modes.
The detector has not yet been descended through the final residual action or
promoted to a relational, causal, asymptotic, or quantum observable.

\section{The normalization triangle}\label{sec:normalization}

Higher-derivative currents are sensitive to integration-by-parts and action
normalization conventions. Three independent constructions are therefore
compared. First, the gauge-reduced axial equation operator is formally
self-adjoint. In Fourier convention $\partial_t\mapsto-i\omega$,
$\partial_x\mapsto ik$, its exact quadratic kernel in the order
$(H_t,H_x,Q_t,Q_x)$ is
\begin{equation}\label{eq:kernel}
K=\begin{pmatrix}
-\frac\lambda4 A&-\frac{k\lambda\omega}{4}B&\lambda&0\\
-\frac{k\lambda\omega}{4}B&\frac\lambda4 C&0&-\lambda\\
\lambda&0&k^2+\lambda&k\omega\\
0&-\lambda&k\omega&-\lambda+\omega^2
\end{pmatrix},
\end{equation}
where
\begin{align}
A={}&3k^4+6k^2\lambda-3k^2\omega^2-2k^2
 +3\lambda^2-3\lambda\omega^2-8\lambda+6\omega^2,\\
B={}&3k^2+3\lambda-3\omega^2+4,\\
C={}&3k^2\lambda-3k^2\omega^2-6k^2+3\lambda^2
 -6\lambda\omega^2-8\lambda+3\omega^4+2\omega^2.
\end{align}
It satisfies $K(\omega,k)^\dagger=K(-\omega,-k)^T$ and is the mixed Hessian
of
\begin{equation}
 S^{(2)}_{\rm red}=\frac12\int
 \Phi(-\omega,-k)^TK(\omega,k)\Phi(\omega,k).
\end{equation}

Second, the Green current generated by this Hessian obeys the exact local
Green identity. Third, independent direct variation of the four-dimensional
Weyl--Maxwell action gives a Lee--Wald current whose integral over $S^2$
equals $N_{\ell m}$ times that reduced Green current. Complete
independent-frequency matrices were checked at $\ell=2,3,4$; natural $SO(3)$
equivariance and the degree-two spectral bound promote the equality to every
$\ell\geq2$ and every $m$. No current improvement remains after
compact-sphere integration.

\begin{proposition}[Hessian--Green--Lee--Wald triangle]
\label{prop:triangle}
In the generic axial block, the exact equation operator is the Hessian of the
reconstructed quadratic reduced action; its local Green current equals the
direct integrated four-dimensional Lee--Wald current with the normalization
used in Theorem~\ref{thm:axial-target}.
\end{proposition}

Proposition~\ref{prop:triangle} is sufficient for the linear phase-space and
signature theorems: it fixes the operator, local conservation law, direct
four-dimensional current, and overall action sign independently. A literal
second expansion of the unreduced four-dimensional action density would be a
valuable additional audit, but it is not used or claimed here.

\section{Where did the graviton go?}\label{sec:interpretation}

The answer depends on which quotient and which boundary problem is being
asked about.

\paragraph{Before the final residual quotient.}
The usual gravitational radiation is represented by the axial and polar
metric-containing normal modes \eqref{eq:dispersion}, together with their
flux-forced Maxwell dressings. Every one of these standard modes has a
nonzero source Lee--Wald pairing and a nonzero target pullback pairing. The
physical $\ell=1$ quotient, generalized homogeneous modes, and twist pairs
are nondegenerate as well. The graviton-like tangent has therefore not been
removed at this stage.

\paragraph{After a global residual quotient.}
A residual cohomology calculation may further identify states under a global
conformal action on a closed compact universe. If its one-particle sector
vanishes, that statement concerns invariants of the fully reduced global
state space. It does not retroactively make the pre-quotient solutions zero,
nor does it show that local wave packets cannot solve the equations. The
coefficient detectors \eqref{eq:detector} likewise exist before this last
descent and may or may not survive it.

\paragraph{Closed cylinder versus scattering.}
The Cauchy surface here is $S^1\times S^2$ and has no null infinity.
Asymptotically flat gravitational radiation is organized by Bondi news,
symplectic flux through $\mathscr I^\pm$, and charged asymptotic symmetries.
Transformations that are proper gauge on a compact slice can acquire boundary
charges and cease to be gauge in that setting. Consequently a compact
residual-cohomology statement cannot be used as evidence that the ordinary
helicity-$\pm2$ scattering space vanishes.

\paragraph{Einstein is a sector, not a Weyl gauge.}
The exact relation established here is
\begin{equation}
 \Sol_{\rm EM}^{\rm lin,std}
 \subsetneq \Sol_{\rm WM}^{\rm lin,ax+standard}
\end{equation}
on the declared background and harmonic strata. The strictness is witnessed
by \eqref{eq:extra-module}. Boundary conditions may later select an Einstein
sector, but such selection is additional dynamics plus boundary data, not a
local Weyl gauge fixing.

\section{Open gates and theorem boundary}\label{sec:open}

The following are not prerequisites for the scoped axial paper, but they are
required to widen its title-level theorem.
\begin{enumerate}
\item \textbf{Polar extra branch.} Repeat the exact operator/module,
Lee--Wald, radical, sign, and detector analysis in even parity. Only then is
the complete additional compact target parity-classified.
\item \textbf{Final residual descent.} Determine which standard and extra
coefficient observables survive the declared global residual quotient. Until
then physical mode means the local-gauge-reduced compact phase space, not the
final residual state space.
\item \textbf{Literal action-density expansion.} This is an independent
normalization audit. The present theorem instead uses the exact
Hessian--Green--direct-Lee--Wald triangle of Proposition~\ref{prop:triangle}.
\item \textbf{Nonlinear extension.} Linear inclusion does not imply that an
Einstein tangent extends to a Weyl--Maxwell family. Fixed-bundle Taub tests
already show scoped second-order obstructions; those belong to a separate
nonlinear charge-fibre paper.
\item \textbf{Causal boundary selection and scattering.} A future theorem
must specify asymptotically flat function spaces, retarded/advanced
complexes, radiative symplectic flux, residual charges, and whether the
Einstein subspace is dynamically closed. Bondi energy, black holes, and
scattering are not inferred here.
\end{enumerate}

The strongest present conclusion is therefore:
\begin{quote}
On the fixed compact product, ordinary Einstein--Maxwell waves embed exactly
and nondegenerately in the linear Weyl--Maxwell phase space, but with a
different Hamiltonian form. Weyl--Maxwell theory also has a complete generic
axial branch of genuine, nonradical additional solutions. Neither branch has
yet been classified after the final residual quotient or as an asymptotic
quantum particle space.
\end{quote}

\section{Exact certification and reproducibility}\label{sec:repro}

Every theorem in this manuscript is a synthesis of exact machine-readable
certificates. The companion claim map
\cert{paper/10-compact-einstein-maxwell-weyl-phase-space-claim-map.json}
pins the full SHA-256 hashes and fails closed if an input changes.

\begin{table}[ht]
\centering
\scriptsize
\begin{tabular}{p{0.29\textwidth}p{0.62\textwidth}}
\toprule
Result & Certificate role\\
\midrule
\cert{EINSTEIN_MAXWELL_PRODUCT_INCIDENCE}
& Common exact product background and fixed-flux rational fixture.\\
\cert{EINSTEIN_MAXWELL_CHEVRETON_TANGENT}
& Complete linear on-shell solution inclusion.\\
\cert{COMPACT_EM_RADIATIVE_SYMPLECTIC_MATCHING}
& Positive source radiative pairing and direct Lee--Wald match.\\
\cert{EINSTEIN_MAXWELL_WEYL_STANDARD_HARMONIC_SYMPLECTIC_INCLUSION}
& Complete standard image, mixed orthogonality, and nondegenerate pullback.\\
\cert{EINSTEIN_MAXWELL_WEYL_AXIAL_OPERATOR}
& Generic target operator, invariant factors, and extra quotient.\\
\cert{EINSTEIN_MAXWELL_WEYL_AXIAL_LEE_WALD_COMPLETION}
& Direct four-dimensional current and complete axial signature.\\
\cert{EINSTEIN_MAXWELL_WEYL_AXIAL_REDUCED_ACTION_HESSIAN}
& Reduced Hessian and normalization triangle.\\
\cert{EINSTEIN_MAXWELL_WEYL_AXIAL_EXTRA_DETECTOR}
& Conserved extra coefficient observables before residual descent.\\
\bottomrule
\end{tabular}
\caption{Principal exact inputs. The claim map also pins the exceptional and
mixed-block component certificates.}
\end{table}

The scoped verification commands are listed below; the breakable path
formatting is typographical and each line is one shell command.
\begingroup\footnotesize
\par\noindent\path{python3 bridge/einstein_sector/verify_einstein_maxwell_product_incidence.py}
\par\noindent\path{python3 bridge/einstein_sector/verify_einstein_maxwell_chevreton_tangent.py}
\par\noindent\path{python3 bridge/einstein_sector/verify_einstein_maxwell_radiative_symplectic_matching.py}
\par\noindent\path{python3 bridge/einstein_sector/verify_einstein_maxwell_weyl_standard_harmonic_symplectic_inclusion.py}
\par\noindent\path{python3 bridge/einstein_sector/verify_einstein_maxwell_weyl_axial_operator.py}
\par\noindent\path{python3 bridge/einstein_sector/verify_einstein_maxwell_weyl_axial_lee_wald_completion.py}
\par\noindent\path{python3 bridge/einstein_sector/verify_einstein_maxwell_weyl_axial_reduced_action_hessian.py}
\par\noindent\path{python3 bridge/einstein_sector/verify_einstein_maxwell_weyl_axial_extra_detector.py}
\par\noindent\path{python3 bridge/einstein_sector/verify_compact_linear_paper_claim_map.py}
\endgroup
The manuscript itself is checked with two passes of \texttt{pdflatex}.

\section{Conclusion}

The Einstein sector is visible without interpretive guesswork. On a common
compact Einstein--Maxwell/Weyl--Maxwell background, the complete standard
Einstein--Maxwell harmonic tangent injects into the Weyl--Maxwell solution
space and remains nondegenerate under the target Lee--Wald form. The identity
inclusion nevertheless changes that form in every class of blocks: by
branch-dependent weights on regular radiation, a factor four on the physical
$\ell=1$ quotient, a unipotent shear on the homogeneous sector, and a factor
minus two on the twists.

The target is strictly larger. Its generic axial quotient contains two
additional nonradical solution summands, detected by conserved coefficient
observables, and the complete direct-current axial signature is $(3,1)$.
These statements settle where ordinary waves and one genuine extra branch
sit before the final residual quotient. They do not turn a classical sign
into a quantum ghost theorem, and they do not replace the separate causal
problem of selecting an Einstein scattering sector at infinity.

\begin{thebibliography}{99}
\bibitem{Bach}
R.~Bach,
\emph{Zur Weylschen Relativit\"atstheorie und der Weylschen Erweiterung des
Kr\"ummungstensorbegriffs},
Math. Z. \textbf{9} (1921) 110--135.

\bibitem{Chevreton}
M.~Chevreton,
\emph{Sur le tenseur de super\'energie du champ \'electromagn\'etique},
Nuovo Cimento \textbf{34} (1964) 901--913.

\bibitem{LeeWald}
J.~Lee and R.~M.~Wald,
\emph{Local symmetries and constraints},
J. Math. Phys. \textbf{31} (1990) 725--743.

\bibitem{IyerWald}
V.~Iyer and R.~M.~Wald,
\emph{Some properties of Noether charge and a proposal for dynamical black
hole entropy},
Phys. Rev. D \textbf{50} (1994) 846--864,
\href{https://arxiv.org/abs/gr-qc/9403028}{arXiv:gr-qc/9403028}.

\bibitem{Maldacena}
J.~Maldacena,
\emph{Einstein gravity from conformal gravity},
\href{https://arxiv.org/abs/1105.5632}{arXiv:1105.5632}.

\bibitem{DeserJoungWaldron}
S.~Deser, E.~Joung, and A.~Waldron,
\emph{Partial masslessness and conformal gravity},
J. Phys. A \textbf{46} (2013) 214019,
\href{https://arxiv.org/abs/1208.1307}{arXiv:1208.1307}.

\bibitem{Stelle}
K.~S.~Stelle,
\emph{Renormalization of higher-derivative quantum gravity},
Phys. Rev. D \textbf{16} (1977) 953--969.
\end{thebibliography}

\end{document}
